\section{The recordings and the split}
\label{sec:recordings}

This section describes the corpus, how it is divided into recordings used for
training and recordings held back, and how much of it the reported numbers
actually rest on. Figure~\ref{fig:contact} shows the corpus and the division
visually; Section~\ref{sec:results} then uses that division throughout.

The corpus is a single task recorded in one session: a garment, already
flattened, is folded once by two arms under teleoperation. It holds
\num{65} recordings and \num{30141} frames at \SI{30}{\hertz}, giving a little
over sixteen minutes of demonstration. Every frame carries five camera images
and the commanded joint positions of both arms. Four of the five cameras are
tactile, so most of the visual input is gel imagery that changes very little
until the grippers close on cloth.

\begin{figure}[t]
  \centering
  \includegraphics[width=\textwidth]{recordings.png}
  \caption{Every recording in the corpus, one frame from the middle of each,
  ordered as they were recorded. The seven outlined in the second colour are
  held back from training. They are the last seven, not a random seven, which
  is a property of the trainer rather than a choice made here and is what
  Section~\ref{sec:threats} returns to.}
  \label{fig:contact}
\end{figure}

The division is made by the trainer and not by hand, so that the recordings
held back are exactly those the trainer withheld. Asked to hold back a tenth,
it takes the \emph{last} tenth of the recordings of each task, rounded upward:
here, \num{58} recordings for training and the final \num{7} held back. This
matters more than the proportion does. Anything that drifted over the course of
the session --- the lighting, where the garment was placed, wear on the gel ---
sits entirely inside the held-back set. A difference between the two halves is
therefore a difference in time as well as a difference in familiarity, and no
number in this report can separate the two. Section~\ref{sec:threats} says what
would.

The reported errors rest on a sample of those frames rather than all of them,
because each scored frame requires a full forward pass and a comparison against
the recorded future. Sampling one frame in twenty gives \num{1127} scored
frames from the training recordings and \num{111} from the held-back ones. The
second number is small, and it bounds what can be claimed: an effect of a few
per cent on \num{111} frames is not distinguishable from sampling noise, while
a difference of several fold is. Section~\ref{sec:findings} marks which of its
claims fall on each side of that line.
